\section{Уроки 4--5: полёт по прямой и движение по локальным точкам}

\subsection{Цели уроков}
\begin{itemize}
\item Освоить управление через скорости \texttt{set\_manual\_speed}.
\item Освоить навигацию по точкам \texttt{go\_to\_local\_point}.
\item Научиться ждать завершение этапа через \texttt{point\_reached()}.
\end{itemize}

\subsection{Ключевые команды}
\begin{itemize}
\item \texttt{set\_manual\_speed(vx, vy, vz, yaw\_rate)}
\item \texttt{go\_to\_local\_point(x, y, z, yaw)}
\item \texttt{point\_reached()}
\end{itemize}

\subsection{Опорные листинги}
\lstinputlisting[style=GeoskanStyle,language=pythoncustom,caption={Прямолинейное движение по скорости}]{../code/python/examples/mission_examples/set_manual_speed.py}
\lstinputlisting[style=GeoskanStyle,language=pythoncustom,caption={Переходы по локальным точкам}]{../code/python/examples/mission_examples/go_to_local_point.py}

\begin{figure}[H]
\centering
\begin{tikzpicture}[scale=1.0]
\draw[->,thick] (-0.2,0) -- (3.5,0) node[right] {$x$};
\draw[->,thick] (0,-0.2) -- (0,3.0) node[above] {$y$};
\filldraw[GeoskanBlue] (0,0) circle (2pt) node[below left] {$P_0$};
\filldraw[GeoskanBlue] (0,1.8) circle (2pt) node[left] {$P_1$};
\filldraw[GeoskanBlue] (1.8,1.8) circle (2pt) node[above] {$P_2$};
\draw[GeoskanOrange,thick,-{Latex[length=2mm]}] (0,0) -- (0,1.8);
\draw[GeoskanOrange,thick,-{Latex[length=2mm]}] (0,1.8) -- (1.8,1.8);
\end{tikzpicture}
\caption{Схема маршрута из \texttt{go\_to\_local\_point.py}}
\end{figure}

\begin{taskbox}[Минимальный набор упражнений]
\begin{enumerate}
\item Измените длительность первого участка в \texttt{set\_manual\_speed.py} с 2 до 1 секунды.
\item Добавьте третью точку в \texttt{go\_to\_local\_point.py}.
\item Сделайте разворот на 90$^\circ$ в последней точке.
\end{enumerate}
\end{taskbox}

\subsection{Критерии успешности}
\begin{enumerate}
\item Прямолинейный участок выполняется без заметного дрейфа.
\item Каждая точка подтверждается \texttt{point\_reached()}.
\item Посадка выполняется в пределах безопасной зоны.
\end{enumerate}
